\section{Related Work}\label{sec:related}

\subsection{AI-assisted and agentic software development}
Research on how large language models support programming first focused on interactive code-completion tools. Barke and colleagues identified two recurring modes of interaction with such tools: developers use them either to speed up work whose solution they already understand, or to explore possible solutions when the next step is uncertain \cite{barke2023grounded}. In a controlled user study, Vaithilingam et al. found that GitHub Copilot did not consistently improve task completion time or success, although participants valued it as a starting point. They noted that understanding, editing, and debugging the generated snippets remained a major obstacle \cite{vaithilingam2022expectation}. Evidence on productivity is therefore mixed and depends on context. Peng et al. reported a significant speed-up in a narrow implementation task \cite{peng2023impact}, while Becker et al. found that AI tools from early 2025 increased task completion time by 19\% for experienced open-source developers working on tasks in mature projects they knew well \cite{becker2025impact}. Together, these studies show that the effects of AI support cannot be predicted from model capability alone - they also depend on the user, the task, the codebase, and the effort spent on verification.

Coding agents extend this model of interaction from local suggestions to actions at the level of the whole repository and complete pull requests. Watanabe et al. analyzed 567 Claude Code pull requests in 157 open-source projects and found that 83.8\% of them were eventually merged; among the merged contributions, 45.1\% still needed corrections from a human \cite{watanabe2025agentic}. Li et al. then introduced AIDev, a much larger dataset covering pull requests created by several coding agents, enabling population-level studies of agentic software development on GitHub \cite{li2026aidev}. Studies based on AIDev started to analyze not only the generated changes, but also the human workflows around them. Gao et al. observed that pull requests co-authored by AI were often submitted by contributors with no prior involvement in the repository and were often merged with little explicit feedback \cite{gao2026autopilot}. Peralta et al. further showed that the merge/reject labels alone are not a sufficient measure of agent performance: only 35.7\% of rejected Agentic-PRs in their manually reviewed sample were clear agent errors, while the remaining cases were caused by workflow constraints and a lack of documented reasons for the decision \cite{peralta2026decisions}. These results motivate a view that takes interaction into account, in which acceptance and review discussion describe different aspects of an agent-related contribution.

\subsection{Pull-request acceptance and code-review discussion}
Interpreting pull-request outcomes has long been studied independently of AI. Gousios et al. showed that pull-request decisions and their processing time are linked to a relatively small set of factors at the level of the contribution, the project, and the contributor \cite{gousios2014pull}. Tsay et al. showed that reviewers also use visible social information and prior relationships, alongside the technical features of the contribution \cite{tsay2014influence}. More recent ecosystem-level evidence also indicates that experience gained through prior contributions, issue participation, and collaboration can improve pull-request acceptance, even for contributors who are new to a given project \cite{meijer2025ecosystem}. As a result, acceptance of Agentic-PRs cannot be treated as a direct measure of either the internal quality of the generated code or the developer's ability: it is an outcome shaped jointly by the change, its author, the repository, and its maintainers.

Review comments require similar caution. Modern code review serves several functions beyond finding defects, including knowledge transfer, building awareness, discussing alternatives, and enforcing project conventions \cite{bacchelli2013expectations}. Developers' perception of review quality also depends on the usefulness and accuracy of the feedback, the reviewer's familiarity with the change, and the characteristics of the submitted code \cite{kononenko2016quality}. A higher number of comments may therefore point to needed corrections, but it may also reflect a complex task, an active review culture, explanations, or design discussion. Conversely, a pull request with no comments may have been simple, poorly reviewed, or handled through interactions not visible in the recorded discussion.

\subsection{Developer experience in studies of agentic contributions}

Developer experience is a latent construct, not a field in GitHub. Exploratory repository studies represented it using signals such as account age, prior commits, accepted pull requests, project involvement, technical skills, and collaboration history. Dey et al., for example, modeled developer expertise within an ecosystem-wide ``Skill Space'', rather than reducing it to activity in a single repository \cite{dey2021expertise}. This line of research supports multidimensional representations, but it also highlights an important limitation of validity: observable activity on the platform is only a proxy for experience and may miss professional work done outside public repositories.

The work closest to ours is by Asdaque et al., who also studied AIDev through the lens of developer experience \cite{asdaque2026novice}. They examined 22\,953 Agentic-PRs submitted by 1719 developers and split the users into two groups based on the number of commits over the account's lifetime, normalized by account age. Their lower-experience group submitted larger changes, received 4.52 times more review comments, had a 31\% lower acceptance rate, and waited 5.16 times longer for a decision. This is direct evidence that developer characteristics are linked to the outcomes of agent-supported contributions. At the same time, their binary, single-indicator split leaves open the question of whether this pattern holds with a broader representation of experience and in a larger population of developers.

Our study complements this work rather than repeating it. We combine AIDev with several GitHub activity indicators and identify three relative profiles through unsupervised clustering. We also study agent-use intensity, which is not captured by comparing only the features of the resulting pull requests. This design lets us check whether the relationships reported under a binary, commit-rate-based definition still hold when experience is represented as a multidimensional activity profile.
